\documentclass{article}
\usepackage[UTF8]{ctex}
\usepackage{amsmath,amssymb,amsfonts}
\usepackage{enumitem}
\usepackage{geometry}
\geometry{a4paper, margin=1in}
% === 错题本排版优化（不参与题目标识） ===
\setlength{\parindent}{0pt}
\setlength{\parskip}{0.35em plus 0.10em minus 0.08em}
\setlength{\emergencystretch}{3em}
\setlist{itemsep=0.25em, topsep=0.40em, parsep=0pt}
\allowdisplaybreaks[2]
% === 错题本排版优化结束 ===
\begin{document}
\section*{福州大学2024年高等代数真题}
\section*{108．福州大学2024年高等代数试卷（回忆版）}
一．填空题。\\

\subsection{题目1}
1．不记得具体题目，但记得可以化简成 \(t^{3}+1 \neq 0\) ，题目问的是 \(t=\) \(\_\_\_\_\)．

\subsection{题目2}
2．\(A=\left(\begin{array}{lll}1 & 2 & 3 \\ 2 & a & 6 \\ 3 & 6 & 9\end{array}\right), r(B)=2\) ，且 \(A B=O\) ，则 \(a=\) \(\_\_\_\_\)

\subsection{题目3}
3．\(A, B\) 均为可逆矩阵，\(\left(\begin{array}{cc}O & A \\ B & O\end{array}\right)^{-1}=\)

\subsection{题目4}
4．\(A\) 为 \(n\) 阶矩阵，\(r(A)=m\) ，线性变换 \(\varphi\) 在基 \(\xi_{1}, \xi_{2}, \cdots, \xi_{n}\) 下的矩阵为 \(A\) ，且 \(A X=0\) 的解空间维数为 \(m\) ，则 \(\operatorname{dim}(\operatorname{Im} \varphi)=\) \(\_\_\_\_\)。

\subsection{题目5}
5．\(f\left(x_{1}, x_{2}, \cdots, x_{2 n}\right)=\sum_{i=1}^{n} x_{i} x_{2 n+1-i}\) ，二次型的符号差为 \(\_\_\_\_\)．

二．简答题。\\

\subsection{题目6}
6．已知 \(a_{1 k}=1(k=1,2, \cdots, 2024), a_{2 s}=1(s=1,2, \cdots, 1012), a_{2 m}=-1(m=1013, \cdots, 2024)\) ，且 \(|A|=2024, A\) 为2024阶方阵，求 \(\sum_{k=1}^{1012} A_{1 k}, \sum_{s=1013}^{2024} A_{1 s}\) 。

\subsection{题目7}
7．\(A_{3 \times 2}, B_{2 \times 3}, A B\) 的特征多项式为 \(\lambda^{3}-3 \lambda^{2}+2 \lambda\) ，求 \(B A\) 的特征多项式。

\subsection{题目8}
8．\(e_{1}, e_{2}\) 为空间的一组基，线性变换 \(\varphi, \psi\) 满足

\[
\varphi\left(e_{1}\right)=\beta_{1}, \varphi\left(e_{2}\right)=\beta_{2}, \psi\left(e_{1}+e_{2}\right)=\beta_{1}+\beta_{2}, \psi\left(e_{1}-e_{2}\right)=\beta_{1}-\beta_{2} .
\]

证明：\(\varphi=\psi\) ．

\subsection{题目9}
9．忘了．

\subsection{题目10}
10．\(A\) 是一个给定的三阶方阵（具体数值忘了），求 \(A\) 的不变因子，初等因子，Jordan 标准形．
\end{document}
